\section*{Capítulo \ref{ch:g10:earth-energy-budget} --- El balance energético de la Tierra}

\begin{solution}{exo:g10:earth-energy-budget:1}
La \omterm{def:g10:earth-energy-budget:insolation}{insolación} de un lugar es la potencia de la luz solar que recibe un
metro cuadrado horizontal en la parte alta de la atmósfera situada sobre él
(promediada sobre un día o un año). El \omterm{def:g10:earth-energy-budget:albedo}{albedo} de una superficie es la
fracción de la luz solar que cae sobre ella y que refleja. La nieve reciente
(de 0.8 a 0.9) refleja mucho más que el océano abierto (0.06).
\end{solution}

\begin{solution}{exo:g10:earth-energy-budget:2}
Reflejada: $100/340 \approx \qty{29}{\percent}$; absorbida por el aire:
$80/340 \approx \qty{24}{\percent}$; absorbida por la superficie: $160/340
\approx \qty{47}{\percent}$.
\end{solution}

\begin{solution}{exo:g10:earth-energy-budget:3}
Asfalto: $800 \times (1 - 0.04) = \qty{768}{W}$. Nieve: $800 \times (1 -
0.85) = \qty{120}{W}$, más de seis veces menos.
\end{solution}

\begin{solution}{exo:g10:earth-energy-budget:4}
La evaporación: \qty{82}{W/m^2}, frente a \qty{21}{W/m^2} del calentamiento
del aire. El agua que se evapora toma su \omterm{def:g10:earth-energy-budget:latent-heat}{calor latente} del mar; el vapor sube
con el aire y, donde se condensa en gotitas de nube, a menudo kilómetros más
arriba, devuelve esa energía al aire que lo rodea.
\end{solution}

\begin{solution}{exo:g10:earth-energy-budget:5}
En diciembre, el hemisferio norte está inclinado en sentido contrario al
Sol: en el polo, el Sol permanece bajo el horizonte todo el día, así que no
llega luz solar. En el solsticio de diciembre ocurre así en todos los lugares
situados más allá de $90^\circ - 23.44^\circ =
\qty{66.56}{\degree}$ N.
\end{solution}

\begin{solution}{exo:g10:earth-energy-budget:6}
$1361 \sin 90^\circ = \qty{1361}{W}$; $1361 \sin 30^\circ =
\qty{680.5}{W}$; $1361 \sin 10^\circ \approx \qty{236}{W}$. La mitad de $S_0$
requiere $\sin h = 0.5$, es decir, $h = \qty{30}{\degree}$.
\end{solution}

\begin{solution}{exo:g10:earth-energy-budget:7}
Reflejado 100, absorbido por el aire 80, absorbido por la superficie 160:
$100 + 80 + 160 = 340$, toda la luz solar está contabilizada. En la
superficie, $185 = 25 + 160$: lo que no se refleja se absorbe.
\end{solution}

\begin{solution}{exo:g10:earth-energy-budget:8}
Disco: $\pi \times (6.371\times10^6)^2 \approx \qty{1.275e14}{m^2}$; potencia:
$1361 \times 1.275\times10^{14} \approx \qty{1.74e17}{W}$. Esfera: $4\pi
R^2 \approx \qty{5.10e14}{m^2}$. Media: $1.74\times10^{17}/5.10\times
10^{14} \approx \qty{340}{W/m^2}$, un cuarto de $S_0$.
\end{solution}

\begin{solution}{exo:g10:earth-energy-budget:9}
$0.7 \times 5.10\times10^{14} \times 3.16\times10^7 \approx
\qty{1.1e22}{J}$ en un año.
\end{solution}

\begin{solution}{exo:g10:earth-energy-budget:10}
Las nubes reflejan $100 - 53 = \qty{47}{W/m^2}$ más de luz solar, que por
tanto no se absorbe. Retienen $267 - 239 = \qty{28}{W/m^2}$ de infrarrojo.
En el balance final, la Tierra pierde $47 - 28 = \qty{19}{W/m^2}$: las
nubes, en conjunto, la enfrían.
\end{solution}

\begin{solution}{exo:g10:earth-energy-budget:11}
Desde el 10 de mayo, aproximadamente, hasta finales de julio, es decir, en
mayo, junio y julio. En el ecuador, el Sol pasa por el cenit a mediodía en
los dos equinoccios, en marzo y en septiembre; en los solsticios pasa a
\qty{23.44}{\degree} del cenit, así que el ecuador recibe algo menos en junio
y en diciembre.
\end{solution}

\begin{solution}{exo:g10:earth-energy-budget:12}
Cada segundo: $82/(2.47\times10^6) \approx \qty{3.32e-5}{kg}$. Cada día:
$3.32\times10^{-5} \times 86\,400 \approx \qty{2.9}{kg}$, una capa de unos
\qty{2.9}{mm}. Cada año: $3.32\times10^{-5} \times 3.16\times10^7 \approx
\qty{1050}{kg}$, alrededor de \qty{1.05}{m} de agua.
\end{solution}

\begin{solution}{exo:g10:earth-energy-budget:13}
$S_0 \sin\delta = S_0/\pi$ da $\sin\delta = 1/\pi \approx 0.318$, luego
$\delta \approx \qty{18.6}{\degree}$. Entonces $\sin\lambda = 0.318/\sin 23.44^\circ
\approx 0.318/0.398 \approx 0.800$: $\lambda \approx \qty{53.15}{\degree}$ o
$180^\circ - 53.15^\circ = \qty{126.85}{\degree}$. Entre ambos valores,
$\lambda$ crece \qty{73.7}{\degree}, es decir, $73.7/360$ de año, unos
\num{75} días, de mediados de mayo a finales de julio.
\end{solution}

\begin{solution}{exo:g10:earth-energy-budget:14}
El hielo absorbe $1 - 0.6 = 0.40$ de la luz solar, y el océano $1 - 0.06 =
0.94$. Potencia adicional: $(0.94 - 0.40) \times 300 \times 10^{12} \approx
\qty{1.6e14}{W}$, más que toda la fotosíntesis del planeta. La energía
adicional calienta el agua, que funde más hielo en los bordes de los
témpanos, lo que deja al descubierto más agua oscura: el cambio se alimenta
a sí mismo.
\end{solution}

\begin{solution}{exo:g10:earth-energy-budget:15}
Cuando la Tierra está más cerca del Sol: $1361 \times (149.6/147.1)^2 \approx \qty{1408}{W/m^2}$; cuando está más lejos: $1361 \times (149.6/152.1)^2 \approx \qty{1317}{W/m^2}$; una diferencia de alrededor de un \qty{7}{\percent}. A
\qty{60}{\degree} N, la \omterm{def:g10:earth-energy-budget:insolation}{insolación} es veinte veces mayor en junio que en
diciembre: una variación de unos pocos puntos porcentuales no puede
producirla. Además, la Tierra está más cerca del Sol a principios de enero,
en pleno invierno boreal, y los dos hemisferios, a la misma distancia del
Sol, tienen estaciones opuestas: lo que produce las estaciones es la
inclinación del eje, no la distancia.
\end{solution}

\begin{solution}{pb:g10:earth-energy-budget:1}
\textbf{1.} $\pi \times (6.371\times10^6)^2 \approx \qty{1.275e14}{m^2}$.

\textbf{2.} $1361 \times 1.275\times10^{14} \approx \qty{1.735e17}{W}$.

\textbf{3.} $4\pi R^2 \approx \qty{5.101e14}{m^2}$.

\textbf{4.} $1.735\times10^{17}/5.101\times10^{14} \approx \qty{340}{W/m^2}$,
es decir, $S_0/4$.

\textbf{5.} Ecuador: \qty{1361}{W}; \qty{60}{\degree} N: $1361 \sin
30^\circ = \qty{680.5}{W}$.

\textbf{6.} $100/340 \approx 0.294$.

\textbf{7.} $340 - 100 = \qty{240}{W/m^2}$; sobre toda la Tierra, $240 \times
5.101\times10^{14} \approx \qty{1.224e17}{W}$.

\textbf{8.} $53/340 \approx 0.156$. Las nubes aumentan el \omterm{def:g10:earth-energy-budget:albedo}{albedo} en un factor
$100/53 \approx 1.9$: casi lo duplican.

\textbf{9.} $(240/5.67\times10^{-8})^{1/4} \approx \qty{255}{K}$, unos
\qty{33}{K} por debajo de los \qty{288}{K} de la superficie.

\textbf{10.} $160 + 342 = \qty{502}{W/m^2}$.

\textbf{11.} $82 + 21 + 398 = \qty{501}{W/m^2}$. La superficie gana alrededor
de \qty{1}{W/m^2}: el desequilibrio de \qty{0.7}{W/m^2}, redondeado.

\textbf{12.} $342/160 \approx 2.1$. El aire envía a la superficie, en forma
de infrarrojo, más del doble de energía que la luz solar que esta absorbe:
es el efecto invernadero.

\textbf{13.} $82 \times 3.156\times10^7/(2.47\times10^6) \approx
\qty{1050}{kg}$, una capa de unos \qty{1.05}{m} de espesor.

\textbf{14.} $0.7 \times 5.101\times10^{14} \approx \qty{3.6e14}{W}$.

\textbf{15.} $2805/72.07 \approx \qty{38.9}{kJ}$ por gramo de carbono.

\textbf{16.} $104.9\times10^9$ toneladas son $1.049\times10^{17}$~g;
$1.049\times10^{17} \times 38.9\times10^3 \approx \qty{4.08e21}{J}$.

\textbf{17.} $4.08\times10^{21}/3.156\times10^7 \approx \qty{1.29e14}{W}$.

\textbf{18.} $1.29\times10^{14}/5.101\times10^{14} \approx
\qty{0.25}{W/m^2}$.

\textbf{19.} $0.7/0.25 \approx 2.8$: el desequilibrio de la Tierra es casi
el triple de la potencia que almacena la fotosíntesis.

\textbf{20.} $1.29\times10^{14}/1.224\times10^{17} \approx 0.0011$: la
fotosíntesis almacena \textbf{alrededor del \qty{0.1}{\percent}} de la luz
solar que absorbe la Tierra, una milésima parte.
\end{solution}
